% =========================================================
%  config/commands.tex
%  Perintah khusus modul praktikum.
% =========================================================

% Nilai awal; diisi ulang oleh metadata.tex dan tiap sesi.
\providecommand{\ModulJudul}{Modul Praktikum Gudang Artefak}
\providecommand{\ModulTeknologi}{Python}
\providecommand{\SesiAktif}{Petunjuk}

% --- Kepala sesi ------------------------------------------
% \sesi{nomor}{judul}{sub-cpmk}{perkiraan waktu}
\newcommand{\sesi}[4]{%
  \clearpage
  \renewcommand{\SesiAktif}{Sesi #1}%
  \phantomsection
  \addcontentsline{toc}{section}{Sesi #1\quad #2}%
  \begin{tcolorbox}[enhanced, colback=ubNavy, colframe=ubNavy,
    boxrule=0pt, arc=2pt, left=16pt, right=16pt, top=13pt, bottom=13pt,
    before skip=0pt, after skip=14pt]
    {\sffamily\bfseries\small\color{ubAccent}%
     SESI #1\ \ \textperiodcentered\ \ #4}\\[4pt]
    {\sffamily\bfseries\LARGE\color{white}#2}\\[7pt]
    {\sffamily\small\color{ubHead}Sub-CPMK #3}
  \end{tcolorbox}%
}

% --- Judul bagian di dalam sesi ---------------------------
\newcommand{\bagian}[1]{\subsection*{#1}}

% --- Nama berkas dan menu --------------------------------
\newcommand{\file}[1]{\texttt{\detokenize{#1}}}
\newcommand{\menu}[1]{\textsf{#1}}

% --- Kotak checkpoint -------------------------------------
\newcommand{\checkpoint}[2]{%
  \begin{catatanbox}[title={Checkpoint Sesi #1}]
  Proyek utuh sesi ini tersedia di \file{code/sesi-#1}. #2
  \end{catatanbox}%
}

% --- Rujukan bab ------------------------------------------
\newcommand{\sumberbab}[1]{%
  {\footnotesize\color{ubMuted}Sumber konsep: #1.}%
}
